% ============================================================
% Chapter 2 -- Statistical mechanics, energy-based models, and diffusion
%
% The context chapter. It exists to make one distinction available to the rest
% of the thesis: the exact score, the empirical score, and the fitted score are
% three different objects, and the difficulties of generative diffusion are
% properties of the gaps between them. Chapter 6 onwards adds a fourth.
% ============================================================

\chapter{Statistical Mechanics, Energy-Based Models, and Diffusion}
\label{ch:statmech}

\section{What statistical mechanics is, and why it is here}
\label{sec:sm-what}

Newton's second law is deterministic and complete. Given the positions and
momenta of $N$ particles and the forces between them, the trajectory of the
system is fixed for all time. For $N$ of order $10^{23}$: the initial condition cannot be written down, the equations cannot be
integrated, and the answer, if it were available, would be a list of
$6 \times 10^{23}$ numbers. Yet the gas
in a room has a temperature, a pressure and a heat capacity, and these are
stable, reproducible, and predictable from a handful of parameters.

Statistical mechanics is the resolution of that tension. It abandons the
trajectory and keeps a probability distribution over microstates, and it fixes
that distribution by a variational principle: among all distributions with a
given mean energy, the one realised at equilibrium maximises the entropy
\begin{equation}
  S[p] \;=\; -\sum_{x} p(x) \log p(x),
  \label{eq:sm-entropy}
\end{equation}
which is the unique measure of uncertainty consistent with a short list of
consistency axioms \citep{jaynes1957information}. Carrying out the maximisation
subject to $\sum_x p(x) E(x) = \bar{E}$ produces the Boltzmann distribution
\begin{equation}
  p(x) \;=\; \frac{1}{Z(\beta)}\, e^{-\beta E(x)},
  \qquad
  Z(\beta) \;=\; \sum_{x} e^{-\beta E(x)},
  \label{eq:sm-boltzmann}
\end{equation}
with the inverse temperature $\beta$ appearing as the Lagrange multiplier of
the energy constraint \citep{landau1980statistical, goldstein2002classical}.

\noindent The maximisation is Toolbox~\ref{tb:sm-boltzmann}.

\begin{toolbox}[tb:sm-boltzmann]{From maximum entropy to the Boltzmann distribution}
Maximise \eqref{eq:sm-entropy} subject to normalisation and the energy
constraint using Lagrange multipliers $\lambda$ and $\beta$:
\[
  \mathcal{L}[p] = -\sum_x p(x)\log p(x)
  - \lambda\Big(\sum_x p(x) - 1\Big)
  - \beta\Big(\sum_x p(x) E(x) - \bar E\Big).
\]
Differentiating pointwise in $p(x)$ for each $x$ and setting the result to
zero,
\[
  \frac{\partial \mathcal L}{\partial p(x)} = -\log p(x) - 1 - \lambda - \beta E(x) = 0
  \quad\Longrightarrow\quad
  p(x) = e^{-1-\lambda}\, e^{-\beta E(x)}.
\]
The prefactor does not depend on $x$; call it $1/Z(\beta)$ and fix it by
normalisation, which forces $Z(\beta) = \sum_x e^{-\beta E(x)}$ and recovers
\eqref{eq:sm-boltzmann}. The multiplier $\beta$ is fixed in turn by demanding
the constraint hold, $\bar E = -\partial_\beta \log Z(\beta)$: with the state
space, the energy function and the reference measure fixed, and on the branch
where that relation is invertible, one number --- the mean energy ---
determines the whole distribution. \hfill$\square$
\end{toolbox}

For the canonical family \eqref{eq:sm-boltzmann}, the equilibrium
thermodynamics follows from the normalising constant: the free energy
is $F = -\beta^{-1}\log Z$, and its derivatives give the energy, the entropy and
the response functions. Two features of that construction matter for what
follows.

The first is that the whole of the macroscopic behaviour
is encoded in a single scalar function of the parameters, $\log Z$, which is a
sum over exponentially many configurations and, for broad classes of discrete
interacting models, is $\#\text{P}$-hard to compute \citep{barahona1982computational,jerrum1993polynomial} --- the tractable
exceptions, trees and Gaussians, are exactly the structures the later chapters
exploit.
The second is that competition between energy, which orders, and entropy, which
disorders, can produce \emph{phase transitions}: the equilibrium state can
change character abruptly as a parameter is varied, even though the microscopic
law does not change at all. A system with a symmetry can settle into a state
that lacks it, and which of the equivalent states it settles into is decided by
fluctuations. The canonical example is the Ising model \citep{ising1925beitrag},
whose two-dimensional case \citet{onsager1944crystal} solved exactly, and which
remains the reference picture for a transition driven by the competition between
energy and entropy. Both features will reappear, unchanged in substance, in the
generative models of the next sections.

\section{Energy-based models: the same object, learned}
\label{sec:sm-ebm}

Traffic between statistical mechanics and learning machines runs in both
directions and is old. The Hopfield network stored memories as minima of an
energy \citep{hopfield1982neural}; its capacity was settled by spin-glass methods
\citep{amit1985storing, mezard1987spin, gardner1988space}; and the Boltzmann
machine and its restricted form made the Boltzmann distribution itself the model
\citep{ackley1985learning, smolensky1986information}. The 2024 Nobel Prize in
Physics recognised that lineage \citep{nobel2024physics}. What follows takes
\eqref{eq:sm-boltzmann} and reads it in the opposite direction. Rather than deriving a distribution from a known energy, it
posits a parametric energy $E_\theta$ and fits it so that the induced
distribution matches data \citep{ackley1985learning, lecun2006tutorial}:
\begin{equation}
  p_\theta(x) \;=\; \frac{e^{-E_\theta(x)}}{Z(\theta)},
  \qquad
  Z(\theta) \;=\; \int_{\R^{d}} e^{-E_\theta(x)}\, \dd x .
  \label{eq:sm-ebm}
\end{equation}
The microstate becomes a data point, the energy becomes a learnable scoring
function, and the partition function keeps its role and its difficulty. Fitting
by maximum likelihood gives a gradient with the structure of a competition,
\begin{equation}
  \nabla_\theta\, \E_{x \sim p_{\mathrm{data}}}\!\left[-\log p_\theta(x)\right]
  \;=\;
  \E_{x^{+} \sim p_{\mathrm{data}}}\!\left[\nabla_\theta E_\theta(x^{+})\right]
  \;-\;
  \E_{x^{-} \sim p_{\theta}}\!\left[\nabla_\theta E_\theta(x^{-})\right],
  \label{eq:sm-ebm-grad}
\end{equation}

\noindent The derivation is Toolbox~\ref{tb:sm-ebm-grad}.

\begin{toolbox}[tb:sm-ebm-grad]{The maximum-likelihood gradient of an energy-based model}
From \eqref{eq:sm-ebm}, $\log p_\theta(x) = -E_\theta(x) - \log Z(\theta)$, so
\[
  \nabla_\theta \log p_\theta(x)
  = -\nabla_\theta E_\theta(x) - \nabla_\theta \log Z(\theta).
\]
The second term does not depend on $x$. Differentiating under the integral
sign in $Z(\theta) = \int e^{-E_\theta(x)}\dd x$,
\[
\begin{aligned}
  \nabla_\theta \log Z(\theta)
  &= \frac{1}{Z(\theta)}
     \int_{\R^d} \nabla_\theta e^{-E_\theta(x)}\,\dd x \\
  &= -\int_{\R^d} \nabla_\theta E_\theta(x)\,
     \frac{e^{-E_\theta(x)}}{Z(\theta)}\,\dd x
   = -\E_{x^- \sim p_\theta}
     \big[\nabla_\theta E_\theta(x^-)\big].
\end{aligned}
\]
Substituting back,
\[
  \nabla_\theta \log p_\theta(x) = -\nabla_\theta E_\theta(x)
  + \E_{x^- \sim p_\theta}\big[\nabla_\theta E_\theta(x^-)\big],
\]
and averaging $-\nabla_\theta \log p_\theta(x)$ over $x = x^+ \sim p_{\mathrm{data}}$
gives \eqref{eq:sm-ebm-grad}. \hfill$\square$
\end{toolbox}

in which the first term lowers the energy of observed data and the second
raises the energy of the model's own samples. The second expectation is over
$p_\theta$, which cannot be sampled directly; the negative phase is what makes
maximum-likelihood training of these models generically intractable, and the
standard approach approximates it by Markov chain Monte Carlo, typically
Langevin dynamics, whose stationary measure is
$\propto e^{-E_\theta}$ \citep{hinton2002training, robert2004monte}. This is
expensive, and its cost is the reason the family was, for a long time,
outperformed by models that give up the explicit density in exchange for cheap
sampling \citep{goodfellow2014generative, kingma2014auto,
papamakarios2021normalizing}. \citet{carbone2025hitchhiker} surveys how the
energy-based view relates to the rest of the generative family.

The learning dynamics of these models are themselves a statistical mechanics
problem, and one with a structure worth carrying forward. The spectral dynamics
of the weights \citep{decelle2017spectral, mezard2017meanfield} and the
emergence of compositional features \citep{tubiana2017emergence} are both
statistical-mechanics accounts of training. \citet{bachtis2024cascade} show that
training a restricted Boltzmann machine
proceeds through a \emph{cascade of second-order phase transitions}
(Figure~\ref{fig:sm-bachtis-cascade}): the model
first encodes a coarse, centre-of-mass feature, and then resolves progressively
finer sub-modes, each emergence accompanied by a diverging susceptibility. The
features are learned in order, from the largest scale down. Section
\ref{sec:sm-transitions} finds the same ordering in a different process, and
Chapter~\ref{ch:model} onwards works in a model where the ordering is
analytic.

\begin{figure}[t]\centering
\includegraphics[width=\linewidth]{fig_bachtis_cascade.png}
\caption[Cascade of phase transitions in an energy-based model]{\emph{Left:}
the eigenvalues $\tilde w_\alpha^2(t)$ of the RBM's weight matrix during
training, each crossing $y=1$ at a time ($t_{\mathrm I}$, $t_{\mathrm{II}}$)
that the theory predicts from the initial conditions alone. \emph{Right:} the
free energy in the order-parameter plane $(h_1,h_2)$ at four of those times.
One minimum splits into two along the coarse feature $\eta^1$, then each of
those splits again, giving four --- the diverging susceptibility at each split
is the second-order transition. Reproduced from \citet{bachtis2024cascade},
their Figure~6.}
\label{fig:sm-bachtis-cascade}
\end{figure}

\section{Generative diffusion}
\label{sec:sm-diffusion}

\begin{figure}[t]\centering
\includegraphics[width=\linewidth]{fig_song_dog.png}
\caption[Forward and reverse SDEs, on data]{The picture the equations below
make precise. The forward SDE corrupts a data point into noise; the reverse
SDE run with the score $\nabla_x \log p_t$ undoes exactly that corruption,
step for step. Reproduced from \citet{song2021scorebased}, their Figure~1.}
\label{fig:sm-sde-dog}
\end{figure}

A diffusion model destroys structure in a controlled way and then learns to run
the destruction backwards (Figure~\ref{fig:sm-sde-dog})
\citep{sohl2015deep, song2019generative,
ho2020denoising, song2021scorebased}. Its deterministic counterparts, neural
ODEs and flow matching \citep{chen2018neural, lipman2023flow}, transport the
same measures along a probability-flow trajectory instead of a stochastic one.
The forward process is a stochastic differential equation
\begin{equation}
  \dd X_t \;=\; f(X_t, t)\,\dd t \;+\; g(t)\,\dd W_t ,
  \label{eq:sm-forward}
\end{equation}
started at $X_0 \sim p_0$, the data law, and producing marginals $p_t$ that
approach a simple terminal measure.

\subsection{The forward process is a relaxation to a Boltzmann equilibrium}
\label{sec:sm-relaxation}

The choice of \eqref{eq:sm-forward} used throughout this thesis is not an
arbitrary corruption. Overdamped Langevin dynamics in a potential $E$,
\begin{equation}
  \dd X_t \;=\; -\nabla E(X_t)\,\dd t \;+\; \sqrt{2\beta^{-1}}\;\dd W_t ,
  \label{eq:sm-langevin}
\end{equation}
has, for a confining potential --- $e^{-\beta E}$ normalisable and the
dynamics non-explosive and ergodic --- a unique stationary density proportional
to $e^{-\beta E(x)}$: the
Boltzmann distribution \eqref{eq:sm-boltzmann} of Section~\ref{sec:sm-what},
recovered as the equilibrium of a dynamics rather than as the solution of a
variational problem \citep{risken1989, gardiner2009stochastic}. Setting
$\beta = 1$ and $E(x) = \tfrac12 \|x\|^2$ gives $f(x,t) = -x$,
$g = \sqrt{2}$: the Ornstein--Uhlenbeck channel, whose transition law is
available in closed form,
\begin{equation}
  x \;=\; e^{-t} a \;+\; \sqrt{\Deltat}\,\varepsilon,
  \qquad \varepsilon \sim \Nd(0, I),
  \qquad \Deltat = 1 - e^{-2t},
  \label{eq:sm-ou}
\end{equation}
and whose stationary law is the standard Gaussian
\citep{uhlenbeck1930theory, gardiner2009stochastic}.

So the terminal measure of the forward process is not merely ``simple noise'':
it is the equilibrium state of a quadratic Hamiltonian at unit temperature. The
forward process takes a system prepared far from that equilibrium --- in the
data distribution, which is the thing with all the structure --- and relaxes it
to a reference equilibrium whose free energy is known exactly. That is the
setting of non-equilibrium thermodynamics, and it is where the construction
came from.

\subsection{Why a path, and not a partition function}
\label{sec:sm-nonequilibrium}

Section~\ref{sec:sm-what} exposed the central computational obstruction of
equilibrium statistical mechanics: the Boltzmann law is simple once its
partition function is known, but
\[
  Z=\sum_x e^{-\beta E(x)}
\]
requires a sum over all microscopic configurations. Non-equilibrium statistical
mechanics suggests a different way of comparing probability laws. Rather than
characterising the endpoints through their partition functions, one drives the
system between them and studies the probability of entire \emph{trajectories}.

The physical prototype is the Jarzynski--Crooks theory. If a system is driven
between equilibrium states, the random work $W$ performed along its realised
trajectory contains the equilibrium free-energy difference:
\begin{equation}
  \E\!\left[e^{-\beta(W-\Delta F)}\right]=1.
  \label{eq:sm-jarzynski-physical}
\end{equation}
This identity holds even when the protocol is finite-time and irreversible
\citep{jarzynski1997nonequilibrium,crooks1999entropy}. The important conceptual
point is that an equilibrium quantity, $\Delta F$, can therefore be related to
an average over \emph{paths}, rather than obtained by evaluating the endpoint
partition functions directly. The original diffusion construction of
\citet{sohl2015deep} uses precisely this non-equilibrium viewpoint: define a
tractable forward process that destroys structure, and learn a reverse process
that assigns probability to the same paths in the opposite direction.

\paragraph{Forward and reverse path measures.}

Discretise the forward diffusion into $T$ steps. Since the corruption process
is chosen by us, its transition probabilities are known. Starting from
$x_0\sim p_0$, its joint path law is
\begin{equation}
  Q(x_{0:T})
  =
  p_0(x_0)
  \underbrace{
  \prod_{k=1}^{T}
  q(x_k\mid x_{k-1})
  }_{q(x_{1:T}\mid x_0)}.
  \label{eq:sm-forward-path}
\end{equation}
Generation instead starts from the tractable terminal law $\pi$ introduced in
Section~\ref{sec:sm-relaxation} and moves backwards through learned
conditionals,
\begin{equation}
  P_\theta(x_{0:T})
  =
  \pi(x_T)
  \prod_{k=1}^{T}
  p_\theta(x_{k-1}\mid x_k).
  \label{eq:sm-reverse-path}
\end{equation}
Thus the model is not specified by one global energy that must subsequently be
normalised. It is built as a product of already normalised local transition
probabilities.

The natural quantity comparing the two processes is their pathwise
log-likelihood ratio (defined when the two path measures are mutually
absolutely continuous, so that the ratio below exists; the Gaussian forward and
reverse transitions used here have full support, and every expectation taken of
the ratio is assumed integrable),
\begin{equation}
  \mathcal W_\theta(x_{0:T})
  :=
  \log
  \frac{Q(x_{0:T})}{P_\theta(x_{0:T})}.
  \label{eq:sm-work}
\end{equation}
When the two path measures are those of a physical forward protocol and its
thermodynamic reversal, this log-ratio is the dimensionless dissipated work
appearing in the Crooks and Jarzynski relations. For a learned
$P_\theta$, it is more generally a \emph{path-space mismatch}; the same
algebra nevertheless survives.

Indeed,
\begin{align}
  \E_Q\!\left[e^{-\mathcal W_\theta}\right]
  &=
  \int
  Q(x_{0:T})
  \frac{P_\theta(x_{0:T})}{Q(x_{0:T})}
  \,\dd x_{0:T}
  \nonumber\\
  &=
  \int
  P_\theta(x_{0:T})
  \,\dd x_{0:T}
  =1.
  \label{eq:sm-jarzynski}
\end{align}
This has the Jarzynski form for dimensionless dissipated work. Taking the
log-ratio itself in expectation gives
\begin{equation}
  \E_Q[\mathcal W_\theta]
  =
  \int Q
  \log\frac{Q}{P_\theta}
  =
  \operatorname{KL}(Q\,\|\,P_\theta)
  \geq 0.
  \label{eq:sm-secondlaw}
\end{equation}
The average mismatch therefore vanishes only if the learned reverse path law
coincides with the forward path law.

\paragraph{From paths to the diffusion training objective.}

For learning, the relevant quantity is not the probability of the complete
path but the likelihood assigned by the model to the observed starting point
$x_0$. Marginalising the intermediate states gives
\begin{equation}
  p_\theta(x_0)
  =
  \int
  P_\theta(x_{0:T})
  \,\dd x_{1:T}.
  \label{eq:sm-model-marginal}
\end{equation}
This integral is still difficult in general. The crucial advantage is that we
possess a tractable distribution over the missing path:
$q(x_{1:T}\mid x_0)$. Multiply and divide by it,
\begin{align}
  p_\theta(x_0)
  &=
  \int
  q(x_{1:T}\mid x_0)
  \frac{
    P_\theta(x_{0:T})
  }{
    q(x_{1:T}\mid x_0)
  }
  \,\dd x_{1:T}
  \nonumber\\
  &=
  \E_{q(x_{1:T}\mid x_0)}
  \left[
    \frac{
      P_\theta(x_{0:T})
    }{
      q(x_{1:T}\mid x_0)
    }
  \right].
  \label{eq:sm-importance-path}
\end{align}
Applying Jensen's inequality to the logarithm yields
\begin{equation}
  \boxed{
  \log p_\theta(x_0)
  \geq
  \E_{q(x_{1:T}\mid x_0)}
  \left[
    \log
    \frac{
      P_\theta(x_{0:T})
    }{
      q(x_{1:T}\mid x_0)
    }
  \right].
  }
  \label{eq:sm-elbo}
\end{equation}
This is the evidence lower bound used by the diffusion construction. Expanding
the two path measures makes its tractability explicit:
\begin{align}
  \mathcal L_{\mathrm{ELBO}}
  &=
  \E_{q(x_{1:T}\mid x_0)}
  \Bigg[
    \log\pi(x_T)
    +
    \sum_{k=1}^{T}
    \log p_\theta(x_{k-1}\mid x_k)
    -
    \sum_{k=1}^{T}
    \log q(x_k\mid x_{k-1})
  \Bigg].
  \label{eq:sm-elbo-expanded}
\end{align}
The forward terms are known; learning acts only on the reverse conditionals.

The gap to the true likelihood is itself a KL divergence:
\begin{equation}
  \log p_\theta(x_0)
  -
  \mathcal L_{\mathrm{ELBO}}
  =
  \operatorname{KL}\!\left(
    q(x_{1:T}\mid x_0)
    \,\middle\|\,
    p_\theta(x_{1:T}\mid x_0)
  \right)
  \geq0.
  \label{eq:sm-elbo-gap}
\end{equation}
The thermodynamic analogy can therefore be stated precisely: irreversibility in
the physical problem appears as a discrepancy between forward and reverse path
measures; in the generative problem, the same path-space discrepancy appears as
the slack of the likelihood bound. Training reduces that discrepancy by making
the learned reverse transitions better approximate the reversal of the chosen
forward process.

This is what replacing a partition function by a path buys. There is no global
normalising constant analogous to $Z(\theta)$ to compute: the terminal law
$\pi$ is normalised, each forward and reverse transition is locally normalised,
and the remaining marginalisation over trajectories is handled through the
known forward path measure. The price is that the reverse transition law must
be learned. The next question is therefore: \emph{what determines the reverse
dynamics?}

\begin{toolbox}[tb:sm-fokkerplanck]{The Fokker--Planck equation, from Itô's lemma}
For a smooth test function $\varphi$, Itô's lemma applied to
\eqref{eq:sm-forward} gives
\[
  \dd\,\E[\varphi(X_t)]
  =
  \E\!\left[
    f(X_t,t)\cdot\nabla\varphi(X_t)
    +
    \frac12 g(t)^2\Delta\varphi(X_t)
  \right]\dd t.
\]
Writing
\[
  \E[\varphi(X_t)]
  =
  \int \varphi(x)p_t(x)\,\dd x
\]
and differentiating with respect to time,
\[
  \frac{\dd}{\dd t}\E[\varphi(X_t)]
  =
  \int
  \varphi(x)\,
  \partial_t p_t(x)\,
  \dd x.
\]
On the right-hand side, integrate the drift term by parts once and the
diffusion term twice:
\[
  \int
  f(x,t)\cdot\nabla\varphi(x)\,
  p_t(x)\,\dd x
  =
  -
  \int
  \varphi(x)
  \nabla_x\!\cdot\!
  \big[f(x,t)p_t(x)\big]
  \,\dd x,
\]
and
\[
  \int
  \Delta\varphi(x)\,
  p_t(x)\,\dd x
  =
  \int
  \varphi(x)\,
  \Delta_x p_t(x)\,
  \dd x,
\]
assuming the boundary terms vanish. Since the resulting identity holds for
every test function $\varphi$,
\begin{equation}
  \boxed{
  \partial_t p_t(x)
  =
  -\nabla_x\!\cdot\!
  \big[f(x,t)p_t(x)\big]
  +
  \frac12 g(t)^2\Delta_x p_t(x).
  }
  \label{eq:sm-fokkerplanck}
\end{equation}
This is the Fokker--Planck, or Kolmogorov forward, equation
\citep{gardiner2009stochastic}. The SDE describes random trajectories; the
Fokker--Planck equation describes how their probability density evolves.
\hfill$\square$
\end{toolbox}

The forward equation contains no score. The score appears when the same
evolution of probability measures is represented in the opposite direction.

Before reversing the stochastic process, it is useful to notice that
\eqref{eq:sm-fokkerplanck} already admits a deterministic representation.
Using
\begin{equation}
  \Delta_x p_t
  =
  \nabla_x\!\cdot\!
  \left(
    p_t\nabla_x\log p_t
  \right),
\end{equation}
the Fokker--Planck equation becomes
\begin{equation}
  \partial_t p_t
  =
  -\nabla_x\!\cdot
  \left[
    \left(
      f
      -
      \frac12 g^2\nabla_x\log p_t
    \right)p_t
  \right].
  \label{eq:sm-probflow-continuity}
\end{equation}
This is a continuity equation. Consequently the deterministic ODE
\begin{equation}
  \dot X_t
  =
  f(X_t,t)
  -
  \frac12 g(t)^2
  \nabla_x\log p_t(X_t)
  \label{eq:sm-probability-flow}
\end{equation}
has the same one-time marginals $p_t$ as the stochastic forward SDE. This is
the \emph{probability-flow ODE} of \citet{song2021scorebased}: individual
trajectories differ from those of the diffusion, but at every time the
distribution of states is the same.

For the Ornstein--Uhlenbeck process fixed in
Section~\ref{sec:sm-relaxation}, Figure~\ref{fig:sm-corruption} shows this
forward loss of structure explicitly. The clean trajectory is progressively
contracted while independent Gaussian noise is accumulated; generation must
undo this evolution.

\begin{figure}[t]\centering
\includegraphics[width=\linewidth]{fig_forward_corruption.pdf}
\caption[The forward channel applied to a correlated sequence]{One clean
stationary AR(1) trajectory with $\alpha=0.85$, in grey, and the same
trajectory after the Ornstein--Uhlenbeck channel \eqref{eq:sm-ou} at four
diffusion times. The signal contracts by $e^{-t}$ while independent noise of
variance $\Deltat$ accumulates. Correlation between neighbouring sites survives
longer than the original amplitude. Generation runs this evolution from right
to left.}
\label{fig:sm-corruption}
\end{figure}

The stochastic reversal is a classical result. For a state-independent
diffusion coefficient $g(t)$, and under regularity conditions on the forward
process --- smooth, strictly positive marginals with an integrable score ---
\citet{anderson1982reverse} shows that the reverse-time process is
\begin{equation}
  \boxed{
  \dd X_t
  =
  \left[
    f(X_t,t)
    -
    g(t)^2\nabla_x\log p_t(X_t)
  \right]\dd t
  +
  g(t)\,\dd\bar W_t,
  }
  \label{eq:sm-reverse}
\end{equation}
when integrated from $T$ back to $0$, starting from $X_T\sim p_T$. It exactly
reproduces the forward marginals in reverse order;
Figure~\ref{fig:sm-sde-overview} draws the forward SDE, the reverse SDE and
the probability-flow ODE over one family of marginals.

The difference between \eqref{eq:sm-probability-flow} and
\eqref{eq:sm-reverse} is worth noting. The deterministic probability-flow ODE
contains $\tfrac12 g^2$ because its velocity field alone must reproduce the
Fokker--Planck evolution. The reverse SDE contains $g^2$ because it must both
change the drift and reverse the effect of the Brownian diffusion. Both depend
on the same object,
\[
  \boxed{
  s_t(x):=\nabla_x\log p_t(x),
  }
\]
the score of the noised distribution.

\begin{figure}[t]\centering
\includegraphics[width=\linewidth]{fig_song_sde.png}
\caption[Forward and reverse SDEs, and the probability-flow ODE]{The forward
SDE carries the data law $p_0$ to a simple terminal law $p_T$. Given the score
$\nabla_x\log p_t$, the reverse SDE follows the same family of marginals in
reverse. The white curves are trajectories of the deterministic
probability-flow ODE, which has the same one-time marginals $p_t$ despite
having no Brownian term. Reproduced from \citet{song2021scorebased}.}
\label{fig:sm-sde-overview}
\end{figure}

This also clarifies the parameterisation used in practical diffusion models.
In the discrete variance-preserving construction of \citet{ho2020denoising},
\begin{equation}
  x_t
  =
  \sqrt{\bar\alpha_t}\,x_0
  +
  \sqrt{1-\bar\alpha_t}\,\varepsilon,
  \qquad
  \varepsilon\sim\Nd(0,I).
  \label{eq:sm-ddpm-marginal}
\end{equation}
The network is usually trained to predict the injected noise,
$\varepsilon_\theta(x_t,t)$, rather than the score itself. These are only two
parameterisations of the same conditional expectation. At the population
optimum,
\begin{equation}
  \boxed{
  s_t(x_t)
  =
  -\frac{
    \E[\varepsilon\mid x_t]
  }{
    \sqrt{1-\bar\alpha_t}
  },
  \qquad
  s_\theta(x_t,t)
  =
  -\frac{
    \varepsilon_\theta(x_t,t)
  }{
    \sqrt{1-\bar\alpha_t}
  }.
  }
  \label{eq:sm-eps-score}
\end{equation}
Thus DDPM noise prediction, continuous-time score matching, the reverse SDE and
the probability-flow ODE all require the same underlying quantity. The next
section separates three different senses in which that quantity can be meant:
the exact score of the population law, the score of the empirical training
distribution, and the fitted score produced by a finite model.

A book-length treatment of diffusion models and their variants is given by
\citet{lai2025principles}.


\section{Three scores}
\label{sec:sm-three-scores}\label{sec:tweedie}

The word ``score'' is used for three different objects in this literature, and
almost every difficulty in generative diffusion is a property of the gaps
between them. The distinction is drawn sharply in \citet{biroli2023generative}
and organises what follows.

\paragraph{The exact score.} The score of the true noised law,
$\nabla_x \log p_t$ with $p_t = p_0 * \Nd(0,\Deltat I)$ after contraction. It is
the object \eqref{eq:sm-reverse} calls for. It has a form that will be used on
every page of this thesis: writing $p_t$ as the channel convolution and
differentiating under the integral gives
\begin{equation}
  \boxed{\;
  \nabla_x \log p_t(x)
  \;=\; \frac{e^{-t}\,\E[\,a \,|\, x\,] - x}{\Deltat}
  \quad\Longleftrightarrow\quad
  \E[\,a \,|\, x\,]
  = e^{t}\big( x + \Deltat\, \nabla_x \log p_t(x) \big).
  \;}
  \label{eq:tweedie}
\end{equation}

\noindent The derivation is Toolbox~\ref{tb:sm-tweedie}.

\begin{toolbox}[tb:sm-tweedie]{The score--posterior identity, from the channel convolution}

\paragraph{Setup.}
The variance--preserving OU channel at time $t$ is the linear--Gaussian map
\[
  x \;=\; e^{-t}\,a \;+\; \sqrt{\Delta_t}\,\xi,
  \qquad a\sim P_0,\quad \xi\sim\mathcal N(0,I_d),\quad \xi\perp a ,
\]
so that $x\mid a\sim\mathcal N(e^{-t}a,\Delta_t I_d)$ with transition density
$p(x\mid a)=G(x;e^{-t}a,\Delta_t)$, where
$G(x;m,v):=(2\pi v)^{-d/2}\exp\!\big(-\lVert x-m\rVert^2/(2v)\big)$.
Marginalising over the prior,
\begin{equation}
  P_t(x)\;=\;\int_{\mathbb R^d} P_0(a)\,G(x;e^{-t}a,\Delta_t)\,\mathrm da .
  \label{eq:marg}
\end{equation}
We assume only $\int\lVert a\rVert\,P_0(a)\,\mathrm da<\infty$. Since $G>0$, we have
$P_t(x)>0$ for every $x$, and the score $S(x,t):=\nabla_x\log P_t(x)$ is well defined.

\paragraph{Step 1: log--gradient identity.}
For any differentiable $f>0$ one has $\nabla\log f=\nabla f/f$. Applied to $P_t$,
\begin{equation}
  S(x,t)\;=\;\frac{\nabla_x P_t(x)}{P_t(x)} .
  \label{eq:loggrad}
\end{equation}
The whole computation is thus $\nabla_x P_t(x)$ divided by $P_t(x)$.

\paragraph{Step 2: differentiate the marginal.}
In \eqref{eq:marg} the prior $P_0(a)$ does not depend on $x$; only $G$ does. Both $G$
and $\nabla_x G$ are bounded, locally uniformly in $x$, by a fixed Gaussian in $a$ times
a polynomial, which is $P_0$--integrable. Hence differentiation passes under the integral (Leibniz's Integral Rule / Dominated Convergence Theorem (DCT)):
\begin{equation}
  \nabla_x P_t(x)\;=\;\int P_0(a)\,\nabla_x G(x;e^{-t}a,\Delta_t)\,\mathrm da .
  \label{eq:diffunder}
\end{equation}

\paragraph{Step 3: gradient of the Gaussian kernel.}
From $\log G(x;m,\Delta_t)=-\lVert x-m\rVert^2/(2\Delta_t)+\text{const}$,
\begin{equation}
  \nabla_x G(x;e^{-t}a,\Delta_t)
  \;=\;-\,\frac{x-e^{-t}a}{\Delta_t}\;G(x;e^{-t}a,\Delta_t) .
  \label{eq:kernelgrad}
\end{equation}
The gradient of the channel is the channel itself times the displacement from its mean,
$-(x-e^{-t}a)/\Delta_t$.

\paragraph{Step 4: substitute and split.}
Insert \eqref{eq:kernelgrad} into \eqref{eq:diffunder}. The factor $x$ is constant in $a$,
so it pulls out and leaves $\int P_0(a)G\,\mathrm da=P_t(x)$ by \eqref{eq:marg}:
\begin{align}
  \nabla_x P_t(x)
  &= -\frac{1}{\Delta_t}\int (x-e^{-t}a)\,P_0(a)\,G(x;e^{-t}a,\Delta_t)\,\mathrm da
  \nonumber\\[2pt]
  &= -\frac{x}{\Delta_t}\,P_t(x)
     \;+\;\frac{e^{-t}}{\Delta_t}\int a\,P_0(a)\,G(x;e^{-t}a,\Delta_t)\,\mathrm da .
  \label{eq:split}
\end{align}
Both integrals are finite: $G\le(2\pi\Delta_t)^{-d/2}$ and $\int\lVert a\rVert P_0<\infty$.

\paragraph{Step 5: Bayes' theorem.}
The joint density of the clean--noisy pair is $P_0(a)\,G(x;e^{-t}a,\Delta_t)$; dividing by
its marginal $P_t(x)$ gives the posterior
\begin{equation}
  P_t(a\mid x)\;=\;\frac{P_0(a)\,G(x;e^{-t}a,\Delta_t)}{P_t(x)} ,
  \label{eq:posterior}
\end{equation}
which is a probability density in $a$ \emph{precisely because} $P_t(x)$ is its normalising
constant. Therefore
\begin{equation}
  \frac{1}{P_t(x)}\int a\,P_0(a)\,G(x;e^{-t}a,\Delta_t)\,\mathrm da
  \;=\;\int a\,P_t(a\mid x)\,\mathrm da\;=:\;\langle a\rangle_{x,t}.
  \label{eq:condmean}
\end{equation}

\paragraph{Step 6: divide by $P_t(x)$.}
Combining \eqref{eq:loggrad}, \eqref{eq:split} and \eqref{eq:condmean}
(linearity of the integral; $P_t(a\mid x)$ a probability measure, so $x$ integrates out):
\begin{equation}
  \boxed{\;S(x,t)\;=\;-\frac{x}{\Delta_t}\;+\;\frac{e^{-t}}{\Delta_t}\,\langle a\rangle_{x,t}\;}
  \qquad
  \langle a\rangle_{x,t}=\int a\,P_t(a\mid x)\,\mathrm da .
  \label{eq:tweedie-tb}
\end{equation}
Coordinate-wise this is $S_k(x,t)=\big(e^{-t}\langle a_k\rangle_{x,t}-x_k\big)/\Delta_t$,
since $\nabla_x$ and the kernel both factorise over components.

\paragraph{Denoiser and noise forms.}
Solving \eqref{eq:tweedie-tb} for the conditional mean,
\begin{equation}
  \langle a\rangle_{x,t}\;=\;e^{t}\big(x+\Delta_t\,S(x,t)\big),
  \label{eq:denoiser}
\end{equation}
so the score is equivalent to the Bayes--optimal denoiser of the clean signal.
Taking $\langle\cdot\rangle_{x,t}$ of the channel relation $\xi=(x-e^{-t}a)/\sqrt{\Delta_t}$
and using \eqref{eq:denoiser} (with $x$ passing through the conditional expectation
unchanged),
\begin{equation}
  \langle \xi\rangle_{x,t}\;=\;-\sqrt{\Delta_t}\,S(x,t).
  \label{eq:noiseform}
\end{equation}

\paragraph{Tilted--measure form.}
Because $(x-e^{-t}a)^2=e^{-2t}(a-e^{t}x)^2$, set $\tilde x:=e^{t}x$ (the noiseless
back--projection of $x$ to $t=0$). Then
$G(x;e^{-t}a,\Delta_t)=(2\pi\Delta_t)^{-d/2}\exp\!\big(-\tfrac{e^{-2t}}{2\Delta_t}\lVert a-\tilde x\rVert^2\big)$,
and \eqref{eq:posterior}--\eqref{eq:condmean} read
\begin{equation}
  \langle a\rangle_{x,t}
  \;=\;\frac{\displaystyle\int a\,P_0(a)\,
        \exp\!\big(-\tfrac{e^{-2t}}{2\Delta_t}\lVert a-\tilde x\rVert^2\big)\,\mathrm da}
       {\displaystyle\int P_0(a)\,
        \exp\!\big(-\tfrac{e^{-2t}}{2\Delta_t}\lVert a-\tilde x\rVert^2\big)\,\mathrm da}.
  \label{eq:tilted}
\end{equation}
This is $a$ averaged under $P_0$ \emph{tilted} by a quadratic pinning to $\tilde x$ with
stiffness $\kappa_t:=e^{-2t}/\Delta_t$. For large $t$, $\kappa_t\to0$ and
$\langle a\rangle_{x,t}\to\int a\,P_0(a)\,\mathrm da$ (prior mean, no information); for
$t\to0^{+}$, $\Delta_t\simeq2t$ gives $\kappa_t\simeq1/(2t)\to\infty$ and,
for a prior with a continuous positive density near $x$,
$\langle a\rangle_{x,t}\to\tilde x\to x$ (data dominates). The regularity
matters: under a discrete prior --- the empirical measure below being the
relevant case --- the tilted mean concentrates on the nearest atoms instead,
and need not approach an arbitrary fixed $x$.

\paragraph{Potential form.}
Let $P_0(a)=Z_0^{-1}e^{-V_0(a)}$ with $V_0\in C^1$, and let $Q_{x,t}(a)\propto
P_0(a)\,e^{-\kappa_t\lVert a-\tilde x\rVert^2/2}$ be the tilted density of
\eqref{eq:tilted}. The Gaussian factor makes $Q_{x,t}$ and its gradient integrable with
vanishing boundary terms, so $\int\nabla_a Q_{x,t}(a)\,\mathrm da=0$. Expanding the
gradient and dividing by the normaliser,
\[
  \big\langle \nabla_a\log P_0(a)-\kappa_t(a-\tilde x)\big\rangle_{x,t}=0
  \quad\Longrightarrow\quad
  \big\langle -\nabla_a V_0\big\rangle_{x,t}=\kappa_t\big(\langle a\rangle_{x,t}-\tilde x\big).
\]
Multiplying by $e^{t}$ and using $e^{t}\kappa_t=e^{-t}/\Delta_t$, $e^{t}\tilde x=e^{2t}x$,
the right side is exactly $-x/\Delta_t+(e^{-t}/\Delta_t)\langle a\rangle_{x,t}=S(x,t)$ by
\eqref{eq:tweedie-tb}. Hence
\begin{equation}
  \boxed{\;S(x,t)\;=\;e^{t}\,\big\langle -\nabla_a V_0\big\rangle_{x,t}\;}
  \label{eq:potentialform}
\end{equation}
the optimal score is $e^{t}$ times the expected original force $-\nabla_a V_0$ under the
tilted measure centred at $\tilde x=e^{t}x$.
\end{toolbox}

We call \eqref{eq:tweedie} the \emph{score--posterior identity}; in the
empirical-Bayes literature it is Tweedie's formula
\citep{robbins1956empirical, miyasawa1961empirical, efron2011tweedie,
vincent2011connection}. It holds for any $p_0$ whatever, and it says something
stronger than an analogy: computing a score \emph{is} solving a Bayesian
inference problem. The exact score is unavailable in practice for the plain
reason that $p_0$ is unknown.

For \emph{vector-valued} data $a \in \R^{L}$ noised coordinatewise,
\eqref{eq:tweedie} holds componentwise but with the \emph{full} posterior:
\begin{equation}
  \big( \nabla_x \log p_t(x) \big)_k
  \;=\; \frac{e^{-t}\, \E[\,a_k \,|\, x_0, \dots, x_{L-1}\,] - x_k}{\Deltat}.
  \label{eq:tweedie-joint}
\end{equation}
The conditioning is on \emph{every} coordinate. The channel is independent
across frames, but the prior is not, so the posterior mean of frame $k$, and
with it the $k$-th component of the score, depends on all $L$ observations.
Equation~\eqref{eq:tweedie-joint} is the exact sense in which the joint score
of a trajectory differs from a stack of per-frame scores, and the reason this
thesis studies the former.

\paragraph{The empirical score.} Replace $p_0$ by the empirical measure of a
training set $\{a^{(i)}\}_{i=1}^{N}$, which puts mass $1/N$ on each example.
The convolution is then a finite Gaussian mixture and the score is available in
closed form,
\begin{equation}
  \nabla_x \log p^{\mathrm{emp}}_t(x)
  \;=\; \frac{1}{\Deltat}\left(
  e^{-t} \sum_{i=1}^{N} w_i(x)\, a^{(i)} \;-\; x \right),
  \qquad
  w_i(x) \;=\; \frac{e^{-\lVert x - e^{-t} a^{(i)}\rVert^2 / 2\Deltat}}
                    {\sum_{j} e^{-\lVert x - e^{-t} a^{(j)}\rVert^2 / 2\Deltat}},
  \label{eq:sm-empirical}
\end{equation}
a softmax-weighted average of training points. This is exactly computable and
exactly useless as a generative model: as $t \to 0$ the weights concentrate on
the nearest training example and the reverse dynamics deposit the trajectory on
top of it. Integrated exactly, \eqref{eq:sm-reverse} driven by
\eqref{eq:sm-empirical} returns the empirical measure itself; run as any
sampler actually is, from a Gaussian start and at finite step size, it returns
points arbitrarily close to the training set. Either way nothing new is
generated, which is the point.


\paragraph{The fitted score.} A neural network $s_\phi(x,t)$ trained by
denoising score matching \citep{hyvarinen2005estimation, vincent2011connection}
to minimise
\begin{equation}
  \E_{t}\, \E_{a \sim p_{\mathrm{data}}}\, \E_{\varepsilon}
  \left\lVert
    s_\phi\big(e^{-t} a + \sqrt{\Deltat}\,\varepsilon,\; t\big)
    + \frac{\varepsilon}{\sqrt{\Deltat}}
  \right\rVert^{2}.
  \label{eq:sm-dsm}
\end{equation}
Two optima must be kept apart. At the population level --- the expectation
over $p_{\mathrm{data}}$ taken literally --- the minimiser of
\eqref{eq:sm-dsm} over all measurable functions is the exact score of the
noised population law \citep{vincent2011connection}. Conditional on a finite
training set, the empirical objective is minimised by the empirical score
\eqref{eq:sm-empirical}. A trained network reaches neither exactly, and one
proposed mechanism of generalisation, in the controlled settings where it has
been analysed, is precisely this gap: the network is constrained enough to be
unable to represent \eqref{eq:sm-empirical}, and that failure to fit the
empirical optimum is what produces novel samples
\citep{kadkhodaie2024generalization, kamb2024analytic}. On that account the
mechanism of generalisation is an approximation-theoretic property of an
architecture, which is a considerably less comfortable foundation than the
variational principle of Section~\ref{sec:sm-what}.

\begin{figure}[t]\centering
\includegraphics[width=\linewidth]{fig_bonnaire_memorize.png}
\caption[Generalisation and memorisation windows, qualitatively]{The
qualitative picture behind Section~\ref{sec:sm-transitions}. \emph{Left:}
sample quality (inverse, blue) and memorised fraction (red) against training
epoch $\tau$; a generalisation window opens between $\tau_{\mathrm{gen}}$ and
$\tau_{\mathrm{mem}}$ whose width grows with the training-set size $n$ ---
along the training-epoch axis, where the mechanism is early stopping. \emph{Right:} the
corresponding regimes in the (data, parameters) plane. Reproduced from
\citet{bonnaire2025memorize}, their Figure~1.}
\label{fig:sm-memorize-summary}
\end{figure}



\section{Transitions along the generation trajectory}
\label{sec:sm-transitions}

Because generation is a dynamical process, its failures are dated. Two
crossover times structure the reverse trajectory.

\begin{figure}[t]\centering
\includegraphics[width=0.62\linewidth]{fig_biroli_regimes.png}
\caption[The three regimes of the reverse trajectory]{The two crossover
times below, illustrated on a two-dimensional Gaussian-mixture example.
Trajectories destined for different classes are indistinguishable above the
speciation time $t_{\mathrm S}$ (regime I); between $t_{\mathrm S}$ and the
collapse time $t_{\mathrm C}$ they commit to a class without yet matching any
training point (regime II, where a model genuinely generalises); below
$t_{\mathrm C}$ they are captured by a single training example (regime III).
Reproduced from \citet{biroli2024dynamical}, their Figure~1.}
\label{fig:sm-regimes}
\end{figure}

\paragraph{Speciation.} Run backwards from noise, the trajectory commits to a
broad class of outcomes long before it commits to a specific sample. In the Gaussian-mixture models analysed by
\citet{biroli2024dynamical} and \citet{achilli2026speciation}, for data whose
covariance has top eigenvalue $\lambda$, the commitment happens at
$t_{\mathrm{S}} \sim \tfrac12 \log \lambda$ (in the variance-preserving
convention used here), when the contracted signal along that direction becomes
comparable with the accumulated noise. When the data carry coherent extensive
structure --- a spiked covariance whose top eigenvalue grows with the
dimension --- $t_{\mathrm{S}}$ diverges logarithmically with it. When the data are hierarchical there is not one such time but
a \emph{ladder} of them, one per level of the hierarchy, and features are
resolved coarse-to-fine \citep{sclocchi2025phase, sclocchi2025probing,
favero2025compositional, favero2025bigger}: the same ordering that
\citet{bachtis2024cascade} observe in the training of an energy-based model,
appearing here along the generation trajectory instead.

\paragraph{Collapse and memorisation.} Later in the reverse trajectory a second
transition occurs, and it is the one that limits the whole construction
(Figures~\ref{fig:sm-regimes} and~\ref{fig:sm-memorize-summary}).
\citet{biroli2024dynamical} and \citet{bonnaire2025memorize} show that dynamics
driven by the empirical score are captured by a single training example past a
time $t_{\mathrm{C}}$, and that avoiding this requires a training set
exponentially large in the data dimension. Memorisation need not be a collapse
onto one example: \citet{achilli2024losing} identify a geometric form in which
dimensions are lost without any single training point capturing the trajectory,
and \citet{garnierbrun2026biased} show how the resulting generalisation can be
biased. Solvable models of the learning problem itself
\citep{cui2025solvable} complete the picture. The criterion is entropic: collapse
occurs when the excess entropy of the noised law relative to the terminal
measure falls below $\log N$, so the dataset size needed to postpone it grows
as $e^{\,\text{(excess entropy)}}$.

The two transitions bracket the useful part of the trajectory. Above
$t_{\mathrm{S}}$ nothing has been decided; below $t_{\mathrm{C}}$ the empirical
score reproduces its training data.

\section{Locality, architecture, and message passing}
\label{sec:sm-unet}

If the fitted score is doing the work, the natural question is what the
architecture contributes, and here a second literature meets the first. The
network in almost every image or video diffusion model is a U-Net
(Figure~\ref{fig:sm-unet}) \citep{ronneberger2015unet}: a convolutional encoder that repeatedly
downsamples, a decoder that upsamples back, and long skip connections joining
each encoder stage to the decoder stage of the same resolution. The design is
almost universal and was, for a long time, justified anecdotally.

\begin{figure}[t]\centering
\includegraphics[width=\linewidth]{fig_unet_architecture.png}
\caption[The U-Net architecture]{The architecture behind the reading below.
Resolution halves at each downsampling step on the encoder side (left) and
doubles again on the decoder side (right); the grey arrows are the skip
connections, copying each encoder stage's feature map across to the decoder
stage of matching resolution before its convolutions run. Reproduced from
\citet{ronneberger2015unet}, their Figure~1.}
\label{fig:sm-unet}
\end{figure}

\citet{mei2024unets} gives it a precise reading, and it is the reading this
thesis works inside. For data generated by a tree-structured hierarchical
model, the exact denoiser is computed by belief propagation, which consists of
an upward pass that accumulates evidence from the leaves and a downward pass
that redistributes it, the second reusing the messages formed by the first. A
U-Net has exactly this shape: the encoder is the upward pass, the decoder is
the downward pass, and the long skip connections are the reuse of the
intermediate messages. On this account, for the hierarchical generative models analysed there, a
U-Net is not merely a good inductive bias for denoising but an
\emph{approximate implementation of the exact inference algorithm}, a reading
that yields sample-complexity bounds for learning the denoiser in those
models; nothing in it certifies an arbitrary U-Net on arbitrary data as
message passing. Convolutional classifiers admit
the corresponding reading for the upward pass alone.

Two consequences are useful for the setting considered in this thesis. First,
sum--product is exact on trees (Theorem~\ref{thm:tree-exact}), so for a chain
the posterior marginals can be computed without introducing an additional
approximation at the inference stage. The chain structure makes the problem
particularly tractable: together with the homogeneous transition law, it leads
in the Gaussian case to the closed-form recursions derived in
Chapter~\ref{ch:gaussian-bp}. This provides a simple reference point against
which learned denoisers can be compared.

Second, the same inference formulation offers a way to study how much of the
surrounding sequence is relevant for estimating the score at a given site.
Rather than treating the receptive field only as an architectural hyperparameter,
one can relate it to the decay of the influence of increasingly distant
observations. Chapter~\ref{ch:gaussian-bp} computes this decay rate explicitly
for the chain studied here. A related result on sequences is
\citet{garnierbrun2024transformers}, who find that a transformer trained on
hierarchically generated sequences reconstructs correlations sequentially,
shortest length scale first, and who use exact belief propagation as the
reference against which the network is measured.

\section{Other generative families}
\label{sec:sm-taxonomy}

To \emph{generate} an
object is to draw a sample from the data law $p_0$, not to evaluate $p_0$ at a
point: there is no single best image, sequence, or molecule, only a
distribution that puts more mass on the ones that look right, and generation
is sampling from it. A training set $\{a^{(i)}\}_{i=1}^N$ is, throughout this
thesis (Section~\ref{sec:sm-three-scores}), exactly this: finitely many
independent draws from $p_0$, standing in for a law nobody can write down.
What every family shares, whatever the recipe, is the same computational
strategy for turning that goal into an algorithm: start from a reference law
that is trivial to sample, almost always an isotropic Gaussian, and construct
a map, deterministic or stochastic, carrying samples of the reference law to
samples that look drawn from $p_0$. Diffusion builds that map by running
\eqref{eq:sm-reverse} backwards from the terminal Gaussian; the families below
build it differently.

\emph{Generative adversarial networks} sample by a single push-forward,
$x = G_\theta(z)$ with $z$ simple noise, and train $G_\theta$ against a
discriminator in a minimax game \citep{goodfellow2014generative}: no
likelihood is ever computed, samples can be sharp, and the adversarial
objective is notoriously unstable to train. \emph{Variational autoencoders}
sample by the same kind of push-forward through a stochastic decoder,
$z \sim p(z)$, $x \sim p_\theta(x \mid z)$, and train by maximising a
tractable lower bound on the likelihood \citep{kingma2014auto};
reconstructions are often blurred, an effect usually attributed to the
Gaussian decoder's conditional averaging rather than to the bound itself.
\emph{Normalizing flows} push a simple base measure through an invertible
network, so the change-of-variables formula gives $\log p_\theta$ exactly
\citep{rezende2015variational, papamakarios2021normalizing}, at the price of
every layer needing a tractable Jacobian. \emph{Autoregressive models} sample
one coordinate at a time, $x_i \sim p_\theta(x_i \mid x_{<i})$, factoring the
joint by the chain rule and training each factor by ordinary classification;
transformer-based language models \citep{vaswani2017attention} and, for
images, PixelCNN-style networks \citep{vandenoord2016pixel} both work this
way, with an exact likelihood but a sampling cost sequential in the data
dimension. \emph{Diffusion models}, the subject of this chapter, instead
sample by repeated denoising (Section~\ref{sec:sm-diffusion}); sampling is
sequential in the noise level rather than the dimension, and iterative rather
than one-shot. \citet{ghio2024sampling} places flows, diffusion and
autoregressive networks under one spin-glass analysis and finds the same
first-order phase transition governing sample quality in all three, so the
comparison between them is not merely architectural.